\gddchapter{Interview sequence}


Below is the guideline for the semi-open guided interview. 
The interview starts with easy warm-up questions and progresses towards harder questions as it goes along. 

The answers were translated from polish back to english.

\section{Introduction}
Remind the participant that this is not a test, rather a joint exploration.

\section{Icebreakers}

\begin{QandA}
   \item Now that you’ve tried both versions, how are you feeling? Was there anything that immediately surprised you?
         \begin{answered}
            Mariusz: Well, the second version is more accessible for me, because I’m used to something like that, and the first one was kind of…
            Something new, something like that, nothing else.
            Although no, I once played a game with voice input — it was an RTS where you gave voice commands to units, maybe you know it.
            What was it called?
            I don’t know.
            I really don’t know.
            But the game was fun.

            Gabriel: Was it an old game?
            
            Mariusz: Older, let’s say from around 2010 to now.
            Definitely not older than 2010.
            And it was kind of timeless, because you controlled it with your voice and sent voice commands to units, like helicopters flying in and blowing things up.

            Gabriel: That’s interesting, I’ll note that down, I should find it later.

            Mariusz: You should find it.
            It’s definitely a psychological RTS real-time game where you plug in a microphone and issue voice commands.
            I just don’t remember whether it was in Polish or more likely in English.
            So yeah, I played something like that once.
            You can do that because it gives you available options, but when you play with voice, you actually have to think about it.
            I also don’t like skipping things, so… oh, are you alive? Are you there?

            Gabriel: Yeah, sorry.

            Gabriel: My Garbo is kind of dried out.

            Mariusz: I don’t like skipping things, so when I see available options, I’ll know whether I can come back to something, whether I can do something, or whether I missed something.

            Gabriel: Yeah, I understand, I understand.

            Mariusz: That’s how I feel.

            Gabriel: So you have a greater awareness of possibilities, basically.

            Mariusz: Yeah, and I like that — I like knowing when that option appears.
            It’s like a summary, because in the first version I don’t have to read all the text to understand everything, but in the second I can deal with it because I see ready-made choice options.

            Gabriel: Yeah, you also don’t have to skip the text here, right?

            Mariusz: So in terms of the first version, it’s much better if you want the text to be fully read and immersive.
            Uh-huh.

            Mariusz: But I’m more in favor of the second version.

            Gabriel: Okay, okay.
            Did anything surprise you when you were playing?

            Mariusz: A lot of text.
            \end{answered}

   \item If you had to describe the difference between these two games to a person who hasn't played either, what would you say?
         \begin{answered}
               Inaudible answer
         \end{answered}
         
\end{QandA}




\section{Grand tour questions}
(remember to pause here 3-5 seconds)

\begin{QandA}


   \item From the moment you first used your voice to navigate, what were you thinking as you decided what to say to the system?

         \begin{answered}
            No, I was basically using the simplest words because I know it’s an AI.
            In the sense that I do it so that it understands it.
            I identify with the character, that it’s me.
            And I’m not saying “he goes there” and “I’m going there.”
         \end{answered}

    \item How was it like to use the voice input in the game?
         \begin{answered}
            Mariusz: Weird, it’s not something everyday.
            Especially with some…

            Gabriel: Sitting next to the bed.

            Mariusz: Very close, so I kind of feel judged.

            Gabriel: But I guess… yeah, I understand.

            Mariusz: In general, when someone is watching me while I’m doing something, I just…
            It also feels weird talking to a computer, always, for me.
            So, like… awkward.
         \end{answered}


\end{QandA}




\section{Core comparative questions}

\begin{QandA}
   \item In which version did you feel more 'inside' the story world?
         \begin{answered}
            Mariusz: Let’s say that in the first one, actually with voice, paradoxically, I had to focus more, read everything, and think about what I’m doing and then say it out loud.

            Gabriel: Okay.

            Mariusz: So let’s say the first one.
           
         \end{answered}

   \item How did the 'effort' of thinking of what to say naturally compare to the 'effort' of knowing which keys to press?

         \begin{answered}
            It’s bigger, because in the second version we have ready-made options you can choose from, while in the first version…
            You need to know what you want to say so it doesn’t become a big mess and so it actually makes sense.
            So you have to think through your decisions, instead of picking pre-made ones provided by the scenario.
         \end{answered}

    \item Did the freedom to say 'anything' make you feel more in control, or was it more overwhelming than having a fixed list of keys?

         \begin{answered}
            Mariusz: But I can’t just say anything.

            Gabriel: Yeah, right, you couldn’t, could you?

            Gabriel: But theoretically you could try, right?

            Gabriel: And the fact that you didn’t have that list of options, just had to formulate things yourself—did that make you feel like you had less control, or did it feel more overwhelming compared to a list of comments?

            Mariusz: I had this feeling—you can get the impression that you have more control over the game and what you’re doing, but it’s kind of false, because obviously it can’t work perfectly if it’s not designed that way.
            We still have specific choices, just hidden between the words in the description.
            But yeah, you could say that voice-wise I feel like I have more choices and more control.
         \end{answered}
    \item How would you compare the overall experience of using the two game versions? Which one would you choose and why?

         \begin{answered}
            Mariusz: Do I have to choose which one is better?

            Gabriel: Yes, that’s also your path.

            Mariusz: For me, as an old gamer and a rather comfort-loving guy, I would choose the third option.
            It feels really familiar to me, I don’t know if…

            Gabriel: The third one?

            Mariusz: The second one.
            I played the third one, sorry.
            I’m not scared.
            I would choose the second option with text provided.
            Because text-based games used to have a place in my heart.
            I don’t know, maybe you know that browser Flash game about dragons, where you take care of your own dragon.

            Gabriel: No, no, I don’t know it.

            Mariusz: That was from the early browser gaming era, at the very beginning, and it was a text-based game.
            And you would pick options there.
            That’s what it reminds me of, so I like it.
            : It’s easy and convenient.
            It is.
            It is, that’s my feeling.
            
         \end{answered}
\end{QandA}






\section{Probing}
In this part of the interview look into the sticky moments noted above and ask follow-up questions probing for more detail.
Include lexical reflection (eg. you used the word "clunky" when explaining navigation. Could you tell me more about what was happening
in that moment?)

The probing was done as follow up questions to the pre-existing questions and in the discussion during the last question rather than separate questions in this test.

% \begin{QandA}
%    \item It looks like you really value creativity when it comes to playing games. Would you agree with that?
%    \begin{answered}
%       I would agree with that because
%       If I played, let's say, a lot of games, then they can become kind of dull.
%       So if something is new for me, or it's like, you know, creative or interesting,
%       I value that because I'm pretty worn out by some normal or usual type of games like the second one.
%       It was way more similar to the games that I played in the past, I would say.
%    \end{answered}
% \end{QandA}


\section{Final reflection}

\begin{QandA}
   \item After reflecting on everything today, is there anything else you’d like to add about using your voice to play this game that we haven't talked about?

         \begin{answered}
            Mariusz: Overall, it’s a nice game, a good idea. I like this kind of storytelling games. I hope there will be more romance options. And the only note I have is that there’s a lot of descriptive text, maybe too much describing what’s around, because it is needed. Especially in insignificant locations, because it takes a bit of time. In that “Babułstaurie” it’s nice. Unless there’s something planned in the future in that album, I don’t know, I’m not aware of the plans. So if something is going to be added in the future, then okay, but…

Gabriel: Yeah, but in the version you played…

Mariusz: In the version I played, it’s fine, it’s nice, it builds the world and so on, but it doesn’t add anything, so… yeah. If the world is…

Gabriel: If you can’t interact with it at all, then…

Mariusz: Sometimes I would prefer, if I’m walking through a corridor that’s just a connector between locations, to have three simple sentences saying it’s a corridor and it smells like shit. That’s it. That would be enough for me. And more concrete things—well, not just more, but more clearly defined things you can actually do.

Gabriel: Yeah, yeah, so if something is described, it should somehow connect to gameplay mechanics, right?

Mariusz: Yes, not just descriptions of things you can’t interact with, but descriptions of things I can actually interact with.

Gabriel: So everything that is described is in some way interactive.

Mariusz: Okay. And that’s it. I like the graphics. They remind me of those Polish text-based games. They also had small graphics so it wasn’t just empty. And I think, for the vibe, to complete the immersion, some light audio would be good—just sounds and ambience, like when you’re in that gallery: wind between abandoned shops, dripping water, or sounds when picking up objects. I think that would really complete it nicely. Because playing it in silence…

Gabriel: I fully agree. And the location itself would be nice, that’s true. Alright, thank you very much.
         \end{answered}
\end{QandA}

\section{Closing}
Thank the participant for his time and tell them again how much value his input and time gives me.






